% SEGMENT OLP-0115-S01
% Part: first-order-logic
% Chapter: axiomatic-deduction
% Section: axioms-rules-quantifiers

\documentclass[../../../include/open-logic-section]{subfiles}

\begin{document}

\olfileid{fol}{axd}{qua}

\olsection{அளவையடைகளுக்கான அடிகோள்களும் விதிகளும்}

% SEGMENT OLP-0115-S02
\begin{defn}[அளவையடைகளுக்கான அடிகோள்கள்]
அளவையடைகளை ஆளும் \emph{அடிகோள்கள்} பின்வருவனவற்றின் அனைத்து
நிகழ்வுகளும் ஆகும்:
\begin{align}
\ollabel{ax:q1} & \lforall[x][!B] \lif !B(t), \\
\ollabel{ax:q2} & !B(t) \lif \lexists[x][!B].
\end{align}
இங்கு $t$ எந்த மூடிய உறுப்பாகவும் இருக்கலாம்.
\end{defn}

% SEGMENT OLP-0115-S03
\begin{defn}[அளவையடைகளுக்கான விதிகள்]
\item $!B \lif !A(a)$ ஏற்கெனவே !!{derivation} இல் வந்து, $a$ ஆனது
  $\Gamma$ அல்லது~$!B$ இல் வரவில்லை எனில், $!B \lif
  \lforall[x][!A(x)]$ ஒரு சரியான அனுமானப் படி.
\item $!A(a) \lif !B$ ஏற்கெனவே !!{derivation} இல் வந்து, $a$ ஆனது
  $\Gamma$ அல்லது~$!B$ இல் வரவில்லை எனில், $\lexists[x][!A(x)] \lif
  !B$ ஒரு சரியான அனுமானப் படி.
\end{defn}

% SEGMENT OLP-0115-S04
இவ்விரண்டில் எதனையும் ``\QR'' எனச் சுருக்கி எழுதுவோம்.

\end{document}

